\documentclass[12pt]{article}
\usepackage[T1]{fontenc}
\usepackage[utf8]{inputenc}
\usepackage{lmodern}
\usepackage{textcomp}
\usepackage{enumitem}
\usepackage{geometry}
\geometry{margin=1in}
\begin{document}
\begin{center}
Answer Key - Psychology - Graduate Level Assessment 13 \
\small (Answers are ordered exactly as in the question paper.)
\end{center}

\section*{Multiple Choice}
\begin{enumerate}[label=\arabic*.]
\item \textbf{Answer:} A
\\ \textit{Options:} A) phonemes from the child's native language only.; B) Phonemes from the language most frequently spoken around the child.; C) Phonetic sounds specific to the child's geographical location.; D) Phonemes from the parents' native language only.; E) morphemes that the child has heard most frequently.
\\ \textit{Note:} Correct option: A - phonemes from the child's native language only.
\item \textbf{Answer:} D
\\ \textit{Options:} A) associative; B) distributive; C) additive; D) disjunctive; E) divisive
\\ \textit{Note:} Correct option: D - disjunctive
\item \textbf{Answer:} D
\\ \textit{Options:} A) avoid setting any specific goals to maintain flexibility; B) focus on developing a deep understanding of the crisis before setting goals; C) try to ensure that the client becomes calm and relaxed; D) agree upon clear-cut goals and keep the sessions focused on these goals; E) primarily use medication as a means of managing the crisis
\\ \textit{Note:} Correct option: D - agree upon clear-cut goals and keep the sessions focused on these goals
\item \textbf{Answer:} A
\\ \textit{Options:} A) increase access to mental health services; B) create a support system for people with mental health problems; C) educate the community about mental health issues; D) prevent community disintegration; E) deal with problems before they occur
\\ \textit{Note:} Correct option: A - increase access to mental health services
\item \textbf{Answer:} B
\\ \textit{Options:} A) short-term memory; B) sensory memory; C) eidetic memory; D) episodic memory; E) explicit memory
\\ \textit{Note:} Correct option: B - sensory memory
\end{enumerate}

\section*{True / False}
\begin{enumerate}[label=\arabic*.]
\setcounter{enumi}{5}
\item \textbf{Answer:} True
\\ \textit{Options:} A) True; B) False
\\ \textit{Note:} LLM-generated true statement based on MCQ.
\item \textbf{Answer:} True
\\ \textit{Options:} A) True; B) False
\\ \textit{Note:} LLM-generated true statement based on MCQ.
\item \textbf{Answer:} True
\\ \textit{Options:} A) True; B) False
\\ \textit{Note:} LLM-generated true statement based on MCQ.
\item \textbf{Answer:} False
\\ \textit{Options:} A) True; B) False
\\ \textit{Note:} LLM-generated false statement. Correct answer: 'fundamental attribution bias'.
\item \textbf{Answer:} True
\\ \textit{Options:} A) True; B) False
\\ \textit{Note:} LLM-generated true statement based on MCQ.
\end{enumerate}

\section*{Long Form Response}
\begin{enumerate}[label=\arabic*.]
\setcounter{enumi}{10}
\item \textbf{Answer:} 1. The relative lack of emotion when dealing with people. 2. Lack of conventional morality. 3. No strong ideological commitments. 4. Relative psychological balance.. Explain why this is the correct answer and provide supporting evidence. Reference key facts or concepts that support this choice.
\\ \textit{Note:} Long-form derived from MCQ; award credit for stating the correct idea and giving supporting reasoning.
\item \textbf{Answer:} Susan's academic performance. Explain why this is the correct answer and provide supporting evidence. Reference key facts or concepts that support this choice.
\\ \textit{Note:} Long-form derived from MCQ; award credit for stating the correct idea and giving supporting reasoning.
\end{enumerate}

\vfill
\noindent\textit{End of Answer Key}
\end{document}